\documentclass[12pt]{article}
%%%%%%%%%%%%%%%%%%%%%%%%%%%%%%%%%%%%%%%%%%%%%%%%%%%%%%%%%%%%%%%%%%%%%%%%%%%%%%%%%%%%%%%%%%%%%%%%%%%%%%%%%%%%%%%%%%%%%%%%%%%%%%%%%%%%%%%%%%%%%%%%%%%%%%%%%%%%%%%%%%%%%%%%%%%%%%%%%%%%%%%%%%%%%%%%%%%%%%%%%%%%%%%%%%%%%%%%%%%%%%%%%%%%%%%%%%%%%%%%%%%%%%%%%%%%
\usepackage{amsfonts}
\usepackage{eurosym}
\usepackage{geometry}
\usepackage{amsmath,amsthm,amssymb}
\usepackage{ulem} 
\usepackage{graphicx}
\usepackage{comment}
%\usepackage[sort,comma]{natbib}
\usepackage[utf8]{inputenc}
\usepackage{setspace}
\usepackage[backend=biber, style = apa]{biblatex}
\usepackage{placeins} % to separate sections

\usepackage{adjustbox}
\usepackage{array}
\usepackage{multirow}
\usepackage{graphicx}
\usepackage{subcaption}
\usepackage{pifont}
\usepackage{amssymb}
\usepackage{comment}
\usepackage[hang, flushmargin, bottom]{footmisc}
\usepackage{footnotebackref}
\usepackage{xcolor}
\usepackage{hyperref}
\usepackage{booktabs}
\usepackage{pifont}
\usepackage{caption}
\usepackage{float}
\usepackage{todonotes}
\setcounter{MaxMatrixCols}{10}


%\setlength{\bibsep}{0.3pt}
\setlength{\textfloatsep}{5pt}
\hypersetup{breaklinks=true,hypertexnames=false,colorlinks=true,citecolor = teal}
\captionsetup{font=normalsize}
\newcommand{\cmark}{\ding{51}}
\def\sym#1{\ifmmode^{#1}\else\(^{#1}\)\fi}
\renewcommand{\thetable}{\Roman{table}}
\geometry{verbose,tmargin=.9in,bmargin=1in,lmargin=.8in,rmargin=.8in,nomarginpar}
\makeatletter
\DeclareTextSymbolDefault{\textquotedbl}{T1}
\theoremstyle{plain}
\newtheorem{thm}{\protect\theoremname}
\theoremstyle{plain}
\newtheorem{prop}[thm]{\protect\propositionname}
\theoremstyle{definition}  % Add this line
\newtheorem{definition}[thm]{Definition}  % Add this line
\theoremstyle{remark}  % Add this line
\newtheorem{remark}[thm]{Remark}  % Add this line
\providecommand{\propositionname}{Proposition}
\providecommand{\theoremname}{Theorem}
\makeatother
\newtheorem{ass}[thm]{Assumption}
% \input{tcilatex}
\usepackage{tikz}
\usetikzlibrary{shapes.geometric, arrows, positioning}


\addbibresource{references.bib}

\newcommand{\lawbox}[1]{\noindent\fbox{\parbox{\dimexpr\textwidth-2\fboxsep-2\fboxrule}{#1}}\medskip}

\begin{document}


\begin{itemize}
    \item \textbf{Que es lo que cambia REDEC respecto a la ley 21.130?  Porque mi impresion es que la ley 21.130 agrega como reportantes de la nomina de deudores a emisores de tarjetas de credito no bancarios (\href{https://www.cmfchile.cl/sitio/aplic/serdoc/ver_sgd.php?s567=aef303581a81b5d6b8c7f134b13d05b7VFdwQmVVMVVRVFZOUkUweVRWUm5lVTEzUFQwPQ%3D%3D&secuencia=-1&t=1641949996}{fuente}), como serian las tiendas del retail. 
    Pero reiteradas veces vemos que REDEC busca extender el sistema de reporte a oferentes de creditos no bancarios (OCNB), por ejemplo ver las presentaciones de Marcel (2022, \href{a. https://www.
    camara.cl/verDoc.aspx?prmID=253759&prmTipo=DOCUMENTO COMISION}{fuente}) o de Berstein (2020, \href{https://www.bcentral.cl/documents/33528/133217/sbj27112020.pdf/1b00823d-93f4-193c-14ee-133b703af529?t=1693928559800}{fuente}). Por lo cual no me queda claro a que se refieren con OCNB y si es que incluye a emsisores de tarjetas de credito no bancarios. }


    "La principal ganancia de una eventual extensión del sistema de información en Chile sería acceder a información positiva respecto de los oferentes de crédito no  bancarios" (\cite{berstein_desarrollo_2020})



    "As a result,  banks may learn a borrower’s total bank debt and bank defaults, but may only observe  reported defaults from non-banks (i.e., cannot access non-bank debt balances). In turn,  non-banks can only learn an individuals’ bank and non-bank defaults from the credit  bureau, but not the level of bank or non-bank consumer credit." (\cite{liberman_equilibrium_2018})


    \item \textbf{Dentro de la informacion que incluye el R04 (la estructura de los registros se encuentra en la circular 2266, anexo 1) se encuentran creditos vigentes y morosidades, pero no hay informacion positiva en cuanto a pagos, por ejemplo alguienn que tuvo un credito y ya lo pago. Pero \textcite{berstein_desarrollo_2020} dice que el registro de credito en Chile cuenta con informacion positiva y negativa (\href{ttps : / / www . bcentral . cl / documents / 33528 / 133217 / sbj27112020 . pdf /
    1b00823d-93f4-193c-14ee-133b703af529?t=1693928559800}{fuente}). Entonces no me queda claro que tipo de informacion positiva se compartia a travez del R04. Cita textual "Para el Sistema bancario existe el registro de crédito CMF con información positiva y negativa"  }

    \item \textbf{Relacionado a la pregunta anterior, los reportantes pueden guardar la informacion que consultan del R04 o RDC10 para crear un historial de creditos de los consumidores? o solo pueden consultar la informacion y no guardarla? y si es lo ultimo, como es el enforcement dado que es dificil de monitorear que un banco no guarde la informacion que consulta?}
    
    \item \textbf{Renzo dijo que con el REDEC agrega el archivo RDC10, este remplaza al R04 o es complementario?}
    
    \item \textbf{El informe normativo \textcite{cmf_informe_2021} al referirse a la nomina de deudores dice que para mayor detalle hay que consultar el arhivo D10, ver nota al pie numero 6. Porque es esto? yo pensaba que la nomina de deudores correspondia al R04.  }
    
    El mismo informe dice "Para efectos de enviar la información se utiliza el mismo archivo D10 del Manual del Sistema de Información, considerando los códigos que sean pertinentes; y recibirán el archivo R04, con la información refundida del sistema." Por lo tanto pareciera ser que las empresas comparten el archivo D10 y luego reciben el R04, por lo tanto el D10 es lo que nosotros podriamos tener acceso a y luego tenemos que ser concientes de que el R04 es al cual las empresas tienen acceso. 

\end{itemize}

\vspace{2cm}


Policies que se pueden utilizar 
\begin{itemize}
    \item ley redec
    \item ley 21.130 que agrega emisores de credito no bancarios a la nomina de deudores
    \item creo que hay algo como que las tarjetas de banco de falabella y walmart son compradas por un banco y por lo tanto son incluidas al sistema de informacion de deudores
    \item creo que en algun momento la deuda estudiantil es eliminada del sistema de inforamcion. 
    \item CMF y Walmart vendieron los datos 
\end{itemize}












 


Este sitio \href{https://www.cmfchile.cl/portal/estadisticas/617/w3-propertyvalue-29526.html}{fuente} menciona a la circular N1 al referirse al informe de tarjetas de credito no bancarias. Seria bueno checkear esta circular para ver si estas tenian que reportar incluso antes de la ley 21.130. 





- "actualmente los principales emisores de tarjetas de crédito vinculados a las grandes cadenas de comercio detallista, ...  están constituidos como sociedades de apoyo al giro bancario (SAG) y, por lo tanto, la información de sus deudores ya se encuentra incorporada a la nómina." luego procede a enumerar especificar "Tarjeta Ripley (SAG de Banco Ripley), Tarjetas Cencosud (SAG de Scotiabank en sociedad con Cencosud), Tarjeta CMR Falabella (SAG de Falabella) y Tarjeta Líder (SAG de BCI en sociedad con Walmart)." 

"Actualmente la Comisión dispone de información desagregada a nivel de RUT de todos los emisores de tarjetas de crédito no bancarios fiscalizados, pero exclusivamente para fines de fiscalización y con un enfoque contable (principalmente para  evaluar los modelos de provisiones por riesgo de crédito)." Lo que implica que la informacion de tarjetas de credito ya se encontraba disponible y la CMF probablemente todavia la tenga, solo que no era parte de la nomina de deudores, por lo que no podia ser usada por los bancos para evaluar el riesgo crediticio de los clientes.

Footnote 8 dice cuales serian las tarjetas de credito que se agregarian al sistema "Actualmente existen 8 empresas emisoras de tarjetas de crédito no bancarias fiscalizadas por la Comisión, destacando aquellos vinculados a las empresas de retail La Polar, ABC Din, Hites, Tricot, Corona y Farmacia Cruz Verde. Los otros dos  emisores corresponden a dos empresas de menor tamaño ubicadas la VI y XIV región." Si uno ve el parrafo anterior dice que la comision tiene informacion de entidades no fiscalizadas para fines de fiscalizacion, el lenguaje es no preciso, pero mi impresion es que estas 8 tarjetas estaban reportando a la CMf. 


Para efectos de enviar la información se utiliza el mismo archivo D10 del Manual del Sistema de Información12, considerando los códigos que sean pertinentes; y recibirán el archivo R0413, con la información refundida del sistema.  El archivo D10 incluye: - RUT - NOMBRE - TIPO DEUDOR: DIRECTO-INDIRECTO - TIPO DE CRÉDITO: CONSUMO-CONTINGENTE - MOROSIDAD - MONTO   Cabe señalar que ambos archivos, que actualmente tienen un periodicidad mensual, a partir del 2022 pasarán a ser semanales. Por lo tanto, los ETNB deberán considerar estas nuevas condiciones.

"Se acoge la solicitud de ampliar el plazo de implementación, de manera que las instituciones fiscalizadas puedan realizar los desarrollos tecnológicos requeridos. Para dicho efecto, el primer envío de archivos, en el contexto de marcha blanca, se iniciará la segunda quincena de diciembre 2021, extendiéndose hasta el mes de junio de 2022" lo que implica que hay una marcha blanca que puede ayudar en terminos de tener data pre y post implementacion de la reforma. 

No se entrega informacion historica: "En cuanto a la información que la CMF remite a las instituciones habilitadas por la LGB, esta corresponde a un nómina al cierre de las fechas predefinidas, no a información histórica." 

“En tal sentido, en este archivo no se considera el mero uso de la tarjeta como medio de pago.” ([CMF, 2021, p. 10]) \textcolor{red}{a que se refiere que no se considera el uso de la tarjeta como medio de pago?}
\end{document}
